\usepackage{catchfile,environ}
\usepackage{multicol}
\input{\MyPath Boostan-Letter}
\usetikzlibrary{calc}

\makeatletter% Taken from http://tex.stackexchange.com/q/109619/5764
\def\declarenumlist#1#2#3{%
	\expandafter\edef\csname pgfmath@randomlist@#1\endcsname{#3}%
	\count@\@ne
	\loop
	\expandafter\edef
	\csname pgfmath@randomlist@#1@\the\count@\endcsname
	{\the\count@}
	\ifnum\count@<#3\relax
	\advance\count@\@ne
	\repeat}
\def\prunelist#1{%
	\expandafter\xdef\csname pgfmath@randomlist@#1\endcsname
	{\the\numexpr\csname pgfmath@randomlist@#1\endcsname-1\relax}
	\count@\pgfmath@randomtemp 
	\loop
	\expandafter\global\expandafter\let
	\csname pgfmath@randomlist@#1@\the\count@\expandafter\endcsname
	\csname pgfmath@randomlist@#1@\the\numexpr\count@+1\relax\endcsname
	\ifnum\count@<\csname pgfmath@randomlist@#1\endcsname\relax
	\advance\count@\@ne
	\repeat}

\def\pgfmathrandomitemwithoutreplacement#1#2{%
	\pgfmath@ifundefined{pgfmath@randomlist@#2}{\pgfmath@error{Unknown random list `#2'}}{%
		\edef\pgfmath@randomlistlength{\csname pgfmath@randomlist@#2\endcsname}%
		\ifnum\pgfmath@randomlistlength>0\relax%
		\pgfmathrandominteger{\pgfmath@randomtemp}{1}{\pgfmath@randomlistlength}%
		\def\pgfmath@marshal{\def#1}%
		\expandafter\expandafter\expandafter\pgfmath@marshal\expandafter\expandafter\expandafter{\csname pgfmath@randomlist@#2@\pgfmath@randomtemp\endcsname}%
		% Now prune.
		\c@pgfmath@counta=\pgfmath@randomtemp\relax%
		\c@pgfmath@countb=\c@pgfmath@counta%
		\advance\c@pgfmath@countb by1\relax%
		\pgfmathloop%
		\ifnum\c@pgfmath@counta=\pgfmath@randomlistlength\relax%
		\else%
		\def\pgfmath@marshal{\expandafter\global\expandafter\let\csname pgfmath@randomlist@#2@\the\c@pgfmath@counta\endcsname=}%
		\expandafter\pgfmath@marshal\csname pgfmath@randomlist@#2@\the\c@pgfmath@countb\endcsname%
		\advance\c@pgfmath@counta by1\relax%
		\advance\c@pgfmath@countb by1\relax%
		\repeatpgfmathloop%
		\advance\c@pgfmath@counta by-1\relax%
		\expandafter\xdef\csname pgfmath@randomlist@#2\endcsname{\the\c@pgfmath@counta}%        
		\else%
		\pgfmath@error{Random list `#2' is empty}%
		\fi%        
	}}  
	
	\def\pgfmathrandomlistcopy#1#2{%
		\pgfmath@ifundefined{pgfmath@randomlist@#2}{\pgfmath@error{Unknown random list `#2'}}{%
			\edef\pgfmath@randomlistlength{\csname pgfmath@randomlist@#2\endcsname}%
			\pgfmathloop%
			\ifnum\pgfmathcounter>\pgfmath@randomlistlength\relax%
			\else%
			\def\pgfmath@marshal{\expandafter\global\expandafter\let\csname pgfmath@randomlist@#1@\pgfmathcounter\endcsname=}%
			\expandafter\pgfmath@marshal\csname pgfmath@randomlist@#2@\pgfmathcounter\endcsname%
			\repeatpgfmathloop%
			\expandafter\xdef\csname pgfmath@randomlist@#1\endcsname{\pgfmath@randomlistlength}%
		}%  
	}
	\makeatother
	
	
	% Define how each questionblock should be handled
	\newcounter{questionblock}
	\newcounter{totalquestions}
	\NewEnviron{questionblock}{}%
	
	\newcommand{\randomquestionsfrombank}[2]{%
		\CatchFileDef{\bank}{#1}{}% Read the entire bank of questions into \bank
		\setcounter{totalquestions}{0}% Reset total questions counters  ***
		\RenewEnviron{questionblock}{\stepcounter{totalquestions}}% Count every question  ***
		\bank% Process file  ***
		\declarenumlist{uniquequestionlist}{1}{\thetotalquestions}% list from 1 to totalquestions inclusive.
		\setcounter{totalquestions}{#2}% Start the count-down
		\RenewEnviron{questionblock}{%
			\stepcounter{questionblock}% Next question
			\ifnum\value{questionblock}=\randomquestion 
			\par% Start new paragraph
			\BODY% Print question
			\fi
		}%
		\foreach \uNiQueQ in {1,...,#2} {% Extract #2 random questions
			\setcounter{questionblock}{0}% Start fresh with question block counter
			\pgfmathrandomitem\randomquestion{uniquequestionlist}% Grab random question from list
			\xdef\randomquestion{\randomquestion}% Make random question available globally
			\prunelist{uniquequestionlist}% Remove picked item from list
			\bank% Process file
		}}
		\usepackage{multicol,color}
		\definecolor{mycolor}{rgb}{0,0,0}
		\newcounter{question}
		\newenvironment{mcquestions}{\begin{itemize}}{\end{itemize}}
		\newcommand{\question}[1]{\stepcounter{question}\item[\textcolor{mycolor}{\thequestion.}] #1}
		\newsavebox{\answera}
		\newsavebox{\answerb}
		\newsavebox{\answerc}
		\newsavebox{\answerd}
		\newcommand{\answer}[4]{%
			\pgfmathdeclarerandomlist{some numbers}{{1}{2}{3}{4}}
			\pgfmathrandomlistcopy{tmp}{some numbers}
			\pgfmathrandomitemwithoutreplacement\j{tmp}
			\sbox\answera{\ifcase\j\relax\or#1\or#2\or#3\or#4\fi}%
			\pgfmathrandomitemwithoutreplacement\j{tmp}
			\sbox\answerb{\ifcase\j\relax\or#1\or#2\or#3\or#4\fi}%
			\pgfmathrandomitemwithoutreplacement\j{tmp}
			\sbox\answerc{\ifcase\j\relax\or#1\or#2\or#3\or#4\fi}%
			\pgfmathrandomitemwithoutreplacement\j{tmp}
			\sbox\answerd{\ifcase\j\relax\or#1\or#2\or#3\or#4\fi}%
			\ifdim\wd\answera<0.165\textwidth%
			\ifdim\wd\answerb<0.165\textwidth%
			\ifdim\wd\answerc<0.165\textwidth%
			\ifdim\wd\answerd<0.165\textwidth%
			\begin{multicols}{4}
				\begin{itemize}
					\item[\textcolor{mycolor}{الف)}] \usebox{\answera}
					\item[\textcolor{mycolor}{ب)}] \usebox{\answerb}
					\item[\textcolor{mycolor}{ج)}] \usebox{\answerc}
					\item[\textcolor{mycolor}{د)}] \usebox{\answerd}
				\end{itemize}
			\end{multicols}
			\else%
			\ifdim\wd\answera<0.4\textwidth%
			\ifdim\wd\answerb<0.4\textwidth%
			\ifdim\wd\answerc<0.4\textwidth%
			\ifdim\wd\answerd<0.4\textwidth%
			\begin{multicols}{2}
				\begin{itemize}
					\item[\textcolor{mycolor}{الف)}] \usebox{\answera}
					\item[\textcolor{mycolor}{ج)}] \usebox{\answerc}
					\item[\textcolor{mycolor}{ب)}] \usebox{\answerb}
					\item[\textcolor{mycolor}{د)}] \usebox{\answerd}
				\end{itemize}
			\end{multicols}
			\else
			\begin{itemize}
				\item[\textcolor{mycolor}{الف)}] #1
				\item[\textcolor{mycolor}{ب)}] #2
				\item[\textcolor{mycolor}{ج)}] #3
				\item[\textcolor{mycolor}{د)}] #4
			\end{itemize}
			\fi
			\else
			\begin{itemize}
				\item[\textcolor{mycolor}{الف)}] #1
				\item[\textcolor{mycolor}{ب)}] #2
				\item[\textcolor{mycolor}{ج)}] #3
				\item[\textcolor{mycolor}{د)}] #4
			\end{itemize}
			\fi
			\else
			\begin{itemize}
				\item[\textcolor{mycolor}{الف)}] #1
				\item[\textcolor{mycolor}{ب)}] #2
				\item[\textcolor{mycolor}{ج)}] #3
				\item[\textcolor{mycolor}{د)}] #4
			\end{itemize}
			\fi
			\else
			\begin{itemize}
				\item[\textcolor{mycolor}{الف)}] #1
				\item[\textcolor{mycolor}{ب)}] #2
				\item[\textcolor{mycolor}{ج)}] #3
				\item[\textcolor{mycolor}{د)}] #4
			\end{itemize}
			\fi\fi
			\else
			\ifdim\wd\answera<0.4\textwidth%
			\ifdim\wd\answerb<0.4\textwidth%
			\ifdim\wd\answerc<0.4\textwidth%
			\ifdim\wd\answerd<0.4\textwidth%
			\begin{multicols}{2}
				\begin{itemize}
					\item[\textcolor{mycolor}{الف)}] \usebox{\answera}
					\item[\textcolor{mycolor}{ج)}] \usebox{\answerc}
					\item[\textcolor{mycolor}{ب)}] \usebox{\answerb}
					\item[\textcolor{mycolor}{د)}] \usebox{\answerd}
				\end{itemize}
			\end{multicols}
			\else
			\begin{itemize}
				\item[\textcolor{mycolor}{الف)}] #1
				\item[\textcolor{mycolor}{ب)}] #2
				\item[\textcolor{mycolor}{ج)}] #3
				\item[\textcolor{mycolor}{د)}] #4
			\end{itemize}
			\fi
			\else
			\begin{itemize}
				\item[\textcolor{mycolor}{الف)}] #1
				\item[\textcolor{mycolor}{ب)}] #2
				\item[\textcolor{mycolor}{ج)}] #3
				\item[\textcolor{mycolor}{د)}] #4
			\end{itemize}
			\fi
			\else
			\begin{itemize}
				\item[\textcolor{mycolor}{الف)}] #1
				\item[\textcolor{mycolor}{ب)}] #2
				\item[\textcolor{mycolor}{ج)}] #3
				\item[\textcolor{mycolor}{د)}] #4
			\end{itemize}
			\fi
			\else
			\begin{itemize}
				\item[\textcolor{mycolor}{الف)}] #1
				\item[\textcolor{mycolor}{ب)}] #2
				\item[\textcolor{mycolor}{ج)}] #3
				\item[\textcolor{mycolor}{د)}] #4
			\end{itemize}
			\fi\fi
			\else
			\ifdim\wd\answera<0.4\textwidth%
			\ifdim\wd\answerb<0.4\textwidth%
			\ifdim\wd\answerc<0.4\textwidth%
			\ifdim\wd\answerd<0.4\textwidth%
			\begin{multicols}{2}
				\begin{itemize}
					\item[\textcolor{mycolor}{الف)}] \usebox{\answera}
					\item[\textcolor{mycolor}{ج)}] \usebox{\answerc}
					\item[\textcolor{mycolor}{ب)}] \usebox{\answerb}
					\item[\textcolor{mycolor}{د)}] \usebox{\answerd}
				\end{itemize}
			\end{multicols}
			\else
			\begin{itemize}
				\item[\textcolor{mycolor}{الف)}] #1
				\item[\textcolor{mycolor}{ب)}] #2
				\item[\textcolor{mycolor}{ج)}] #3
				\item[\textcolor{mycolor}{د)}] #4
			\end{itemize}
			\fi
			\else
			\begin{itemize}
				\item[\textcolor{mycolor}{الف)}] #1
				\item[\textcolor{mycolor}{ب)}] #2
				\item[\textcolor{mycolor}{ج)}] #3
				\item[\textcolor{mycolor}{د)}] #4
			\end{itemize}
			\fi
			\else
			\begin{itemize}
				\item[\textcolor{mycolor}{الف)}] #1
				\item[\textcolor{mycolor}{ب)}] #2
				\item[\textcolor{mycolor}{ج)}] #3
				\item[\textcolor{mycolor}{د)}] #4
			\end{itemize}
			\fi
			\else
			\begin{itemize}
				\item[\textcolor{mycolor}{الف)}] #1
				\item[\textcolor{mycolor}{ب)}] #2
				\item[\textcolor{mycolor}{ج)}] #3
				\item[\textcolor{mycolor}{د)}] #4
			\end{itemize}
			\fi\fi
			\else
			\ifdim\wd\answera<0.4\textwidth%
			\ifdim\wd\answerb<0.4\textwidth%
			\ifdim\wd\answerc<0.4\textwidth%
			\ifdim\wd\answerd<0.4\textwidth%
			\begin{multicols}{2}
				\begin{itemize}
					\item[\textcolor{mycolor}{الف)}] \usebox{\answera}
					\item[\textcolor{mycolor}{ج)}] \usebox{\answerc}
					\item[\textcolor{mycolor}{ب)}] \usebox{\answerb}
					\item[\textcolor{mycolor}{د)}] \usebox{\answerd}
				\end{itemize}
			\end{multicols}
			\else
			\begin{itemize}
				\item[\textcolor{mycolor}{الف)}] #1
				\item[\textcolor{mycolor}{ب)}] #2
				\item[\textcolor{mycolor}{ج)}] #3
				\item[\textcolor{mycolor}{د)}] #4
			\end{itemize}
			\fi
			\else
			\begin{itemize}
				\item[\textcolor{mycolor}{الف)}] #1
				\item[\textcolor{mycolor}{ب)}] #2
				\item[\textcolor{mycolor}{ج)}] #3
				\item[\textcolor{mycolor}{د)}] #4
			\end{itemize}
			\fi
			\else
			\begin{itemize}
				\item[\textcolor{mycolor}{الف)}] #1
				\item[\textcolor{mycolor}{ب)}] #2
				\item[\textcolor{mycolor}{ج)}] #3
				\item[\textcolor{mycolor}{د)}] #4
			\end{itemize}
			\fi
			\else
			\begin{itemize}
				\item[\textcolor{mycolor}{الف)}] #1
				\item[\textcolor{mycolor}{ب)}] #2
				\item[\textcolor{mycolor}{ج)}] #3
				\item[\textcolor{mycolor}{د)}] #4
			\end{itemize}
			\fi\fi
		}


%\defglsentryfmt[acronym]{\glsentryname{\glslabel}\ifglsused{\glslabel}{}{ (\lr{\glsentrydesc{\glslabel}})}}

%\defglsentryfmt[english]{\glsgenentryfmt\ifglsused{\glslabel}{}{ (\lr{\glsentryname{\glslabel}})}}

\newcommand\tline[2]{$\underset{\text{#1}}{\text{\underline{#2}}}$}

\geometry{a4paper,top=1.3cm,bottom=1.2cm,left=.9cm,right=.9cm,verbose,nohead,nofoot}

%\newcommand{\insertName}{
%		\baselineskip = .8cm
%		\begin{table}[H]
%			\begin{tabular}{|p{5cm}|p{3cm}|p{3cm}|}\hline
%				نام و نام خانوادگی & شماره دانشجویی & تاریخ 
%				\\\hline
%				& & 
%				\\\hline
%			\end{tabular}
%		\end{table}
%	}


\newcommand{\myabr}[1]{\glsentrytext{#1} (\lr{\glsentrydesc{#1}})}

\NewDocumentEnvironment{bb}{g g}{
	\tikzstyle{mybox1} = [draw=YellowGreen, fill=white,very thick, rectangle, rounded corners, inner sep=10pt, inner ysep=20pt]
	\tikzstyle{fancytitle1} =[fill=YellowGreen, text=white]
	\tikzstyle{fancytitle2} =[fill=YellowGreen!5, text=white]
	\tikzstyle{fancytitle3} =[fill=white, text=white]
	\begin{center}
		\begin{tikzpicture}
			\node [mybox1] (box)\bgroup
			\IfValueTF{#2}{
				\IfFileExists{#2}{\begin{minipage}{.85\textwidth}}{\begin{minipage}{.93\textwidth}}
			}{%%
				\IfFileExists{note.png}{\begin{minipage}{.85\textwidth}}{\begin{minipage}{.93\textwidth}}
			}%%
			\baselineskip=.95cm
				\begin{RTL}
}{%
				\end{RTL}
			\end{minipage}
			\egroup;
			\IfValueTF{#1}{\node[fancytitle1, left=10pt] at (box.north east) {\hboxR{#1}};}{\node[fancytitle1, left=10pt] at (box.north east) {\hboxR{نکته}};}%
			\IfValueTF{#2}{
				\IfFileExists{#2}
				{\node[fancytitle3, left=3pt,   rounded corners] at (box.west) {\includegraphics[width=.07\textwidth]{#2}}; }
				{\node[fancytitle2,  rounded corners] at (box.west) {\handBS};}			
			}{%%
				\IfFileExists{bb.png}
				{\node[fancytitle3, left=3pt,   rounded corners] at (box.west) {\includegraphics[width=.07\textwidth]{note}}; }
				{\node[fancytitle2,  rounded corners] at (box.west) {\handBS};}
			}%%
		\end{tikzpicture}
	\end{center}
}%

\defglsentryfmt[english]{\glsgenentryfmt\ifglsused{\glslabel}{}{}}

\defglsentryfmt[acronym]{\glsentryname{\glslabel}\ifglsused{\glslabel}{}{}}


\newglossarystyle{myAbbrlist2}{%
	%%% این دستور در حقیقت عملیات گروه‌بندی را انجام می‌دهد. بدین صورت که اختصارات‌ در بخش‌های جداگانه گروه‌بندی می‌شوند، 
	%%% عنوان بخش همان نام حرفی است که هر اختصار در آن گروه با آن شروع شده است. 
	\renewenvironment{theglossary}{}{}
	\renewcommand*{\glsgroupskip}{\vskip 0mm}
	\renewcommand*{\glsgroupheading}[1]{}
	%%% در این دستور نحوه نمایش اختصارات می‌آید. در این جا حالت کوچک اختصار در سمت چپ و حالت بزرگ در سمت راست قرار داده شده است، و بین آن با نقطه پر می‌شود. 
	\renewcommand*{\glossentry}[2]{\noindent\glsentrytext{##1}\dotfill\space \Glsentrylong{##1}
		
	}
	%%% تغییر نام محیط abbreviation به فهرست اختصارات
	\renewcommand*{\acronymname}{\rl{فهرست اختصارات}}
}


\newglossarystyle{myEntoFa2}{%
	%%% این دستور در حقیقت عملیات گروه‌بندی را انجام می‌دهد. بدین صورت که واژه‌ها در بخش‌های جداگانه گروه‌بندی می‌شوند، 
	%%% عنوان بخش همان نام حرفی است که هر واژه در آن گروه با آن شروع شده است. 
	\renewenvironment{theglossary}{}{}
	\renewcommand*{\glsgroupskip}{\vskip 0mm}
	\renewcommand*{\glsgroupheading}[1]{}
	%%% در این دستور نحوه نمایش واژه‌ها می‌آید. در این جا واژه فارسی در سمت راست و واژه انگلیسی در سمت چپ قرار داده شده است، و بین آن با نقطه پر می‌شود. 
	\renewcommand*{\glossentry}[2]{\noindent\glsentrytext{##1}\dotfill\space \glsentryname{##1}
		
	}
}

\newtheorem*{definition*}{تعریف}
\newcommand{\putAnswer}[1]{\ifdefined\ShowAnswer\textcolor{Plum}{\textbf{پاسخ: }#1}\fi}
